\section{The first division identity on the common selected realization}
\label{sec:division-identity}

\status{This chapter proves a sufficient conditional constructor for the first, literal \texttt{DivDef3} on the realization fixed in CSR. The proof includes the inner and terminal level identities, the row-projector identity under the explicit row-pair integrability input below, both composite row identities, and the three fields of the first theta subrecord. It retains the original common readings, \texttt{Pins}, \texttt{Holds}, group tie, geometry, rational tests, positive operators, and stored-chain existence. The additional row-pair integrability producer is not proved here, so this is not yet the full original theorem from its original premises alone. The stronger second output \texttt{DivRows3}, including its theta square, is also separate.}

All formal locators refer to the pinned source
\texttt{7cbf7a0c32fca1badf05931879a37c3f8a2b03b4}.
The corresponding declarations are recorded with exact byte ranges in
the accompanying provenance. Inner products are conjugate-linear in
their first argument. An additive ``closure'' in the lattice argument
means the subgroup generated by a set, not topological closure.

% BEGIN DI1
\paragraph{DI1: the common objects and the exact sufficient contract.}
\label{claim:DI1}
Fix the original
\[
 D,G,A_3,K_s,R,F,c_G,X_0,Y,Z,C,s_{EE}.
 \tag{DI.1}
\]
Here \(R:\mathrm{RealDat3}\), \(F:\mathrm{Face3}\),
\(Y:\mathrm{Read}\), \(Z:\mathrm{ReadL}\), and
\(C:\mathrm{ReadC}\) have precisely the dependencies in
\texttt{divTied3\_read}. The coefficient field is
\(L=D.\mathrm{what}_2.LL\), its fixed embedding is
\(\sigma=D.\mathrm{what}_2.\mathrm{embLL}\), and
\(K_s\hookrightarrow L\) is the original inclusion. We retain
\(\mathrm{Pins}(R,F,c_G)\), \(\mathrm{Holds}(R,F)\), and the source's
five-predicate real-group input
\[
 \begin{gathered}
 \mathrm{DiscreteSeriesCover},\quad
 \mathrm{HarishChandraModules},\quad
 \mathrm{CasselmanWallach},\\
 \mathrm{HarishChandraFinite},\quad
 \mathrm{CuspRapidDecay}.
 \end{gathered}
 \tag{DI.2}
\]
These predicates are inputs, not proved by their names. The new level
proof below uses the Harish--Chandra module input. The row metric proof
uses the additional precise integrability input stated next. The
source derives such integrability using its full cusp-decay input
and a growth dictionary whose human reconstruction remains separate.

For \(b\in\{0,1\}\), let \(P_b\) mean the actual level datum
\(R.\mathrm{levR}(F.\mathrm{lv2},b)\), and put
\[
 n_b=\mathrm{nP}(F.\mathrm{lv2},b)>0,\qquad
 \mu_b=n_b^{-1}\nu_0.
 \tag{DI.3}
\]
The latter equality means exactly the source's
\(\mathrm{ENNReal.ofReal}(n_b^{-1})\)-multiple of \(R.\nu_0\),
on the unchanged fundamental domain \(R.\mathrm{fd0}\).
No normalization of \(\nu_0,\nu_Y,\nu_X\) is altered.

We retain the following additional analytic supplier, denoted
\(\mathrm{RowPairInt}\). For every \(b,t\), every \(w\) in
the full original \(P_b\)-level multiplicity space of \(R.MY(0)\),
every \(c\) in its original good part, every pair of smooth tuples
\(v,v'\), and every \(g\in G\), require
\[
 h\longmapsto
 \overline{A_{R.e_Gg}\bigl(c(v)\bigr)(h)}\,w(v')(h)
 \quad\hbox{to be integrable on }\mathrm{fd0}
                         \hbox{ for }\mu_b.
 \tag{DI.3a}
\]
Here \(A\) is the literal normalized finite right translate
(DI.24). Its inverse is included because \(g\) ranges over the
whole finite group. The identity element includes the untranslated
product. Right-translation invariance of the quotient measure gives
the adjoint translate products from the same hypothesis; hence it is
unnecessary to assume a separate integral after applying \(e_0\).
This is a new, expressly retained sufficient input to the human
constructor. It is not being claimed as a consequence of the name
\(\mathrm{CuspRapidDecay}\), or of the special theta-seed bound.

Unpack \(\mathrm{Pins}.\mathrm{char}.\mathrm{base}.\mathrm{chain}\).
Its stored weak chain, together with
\(\mathrm{Holds}.\mathrm{aux}.\mathrm{ex}\), activates the same
\(\omega_Y,\omega_X,v_Y,v_X,\mathrm{ItY},\mathrm{ItX},
\mathrm{ioY},\mathrm{ioX},\mathrm{Th}\).
For each stored existential implication its antecedent is the matching
component of \(\mathrm{ExAt}\); therefore its consequent is true for
the already stored map. This is the complete logical content of
\texttt{chain\_of\_weak}. In particular the resulting
\(\mathrm{IsItW}\) has its original \(\mathrm{ne}\) field, and each
\(\mathrm{IsIo}\) has its original integral equation. SPR25 alone is
not a substitute for this chain input.

The only additional group premise used at the end is the original rule
\[
 \begin{gathered}
 \begin{aligned}
 |G|&=D.\mathrm{what}_2.\mathrm{restRowsFiniteGroupCard},\\
 |F.P_0|&=|F.P_1|(g_0g_1)^{kf}
 \end{aligned}\\
 \Longrightarrow\quad
 \mathrm{GrpTie}(D,A_3.A,F.P_0,F.\psi_0,F.P_1,F.\psi_1,
                       c_G\circ R.e_G).
 \end{gathered}
 \tag{DI.4}
\]
The map in (DI.4) is written in the source order
\(R.e_G.\mathrm{trans}(c_G)\). These hypotheses refer to one
realization, not independently chosen objects with matching dimensions.
% END DI1

% BEGIN DI2
\paragraph{DI2: explicit compact averages approximate every smooth vector.}
\label{claim:DI2}
Let \(H\) be one of the original local Hilbert models, with unitary
action \(\pi\) of \(U(p,q)\), and let \(K=U(p)\times U(q)\) be its
actual maximal compact subgroup in the model's signed frame.
There are finite-rank-on-orbits averaging operators \(A_N\) such that
\(A_Nv\) is \(K\)-finite and, for every smooth \(v\),
\[
             D(A_Nv-v)\longrightarrow0\quad\hbox{in }H
 \tag{DI.5}
\]
for every finite word \(D\) in the derived real Lie algebra operators.
This assertion uses the original smooth vectors and their Sobolev
seminorms. The source's span of smooth vectors consists of smooth
vectors: it takes only finite sums, and differentiation is linear.

\emph{Proof.}
If \(n=p+q=0\), use the identity operator. Otherwise normalize Haar
measure on the compact group \(K\) to mass one and set
\[
 b(k)=\frac{2n+\operatorname{Re}\operatorname{tr}k}{3n},
 \qquad q_N(k)=\frac{b(k)^N}{\int_K b(h)^N\,dh},
 \qquad A_Nv=\int_Kq_N(k)\pi(k)v\,dk.
 \tag{DI.6}
\]
Unitary eigenvalues give \(1/3\le b\le1\), with equality \(b=1\)
only at the identity. The denominator is positive. For a neighborhood
\(U\) of the identity, choose \(c<d<1\) such that \(b\le c\) on
\(K\setminus U\) and \(b>d\) on a smaller neighborhood \(V\).
Then
\[
          \int_{K\setminus U}q_N
              \le \frac{1}{\operatorname{vol}(V)}(c/d)^N
              \longrightarrow0.
 \tag{DI.7}
\]
The polynomial \(b^N\) belongs to the finite-dimensional space of
polynomials of degree at most \(N\) in the matrix entries and their
conjugates, restricted to \(K\). Left translates preserve that space.
Changing variables in (DI.6) therefore places the whole \(K\)-orbit
of \(A_Nv\) in the span of finitely many integrals; hence \(A_Nv\)
is \(K\)-finite.

For a Lie word \(D=X_1\cdots X_a\), differentiating the unitary
orbit gives
\[
 D\pi(k)v=\pi(k)
       d\pi(\operatorname{Ad}(k^{-1})X_1)\cdots
       d\pi(\operatorname{Ad}(k^{-1})X_a)v.
 \tag{DI.8}
\]
For smooth \(v\) this is a continuous, bounded function of \(k\)
on \(K\), with value \(Dv\) at one. Differentiation under the
compact integral is justified by the corresponding uniform bound on
one further derivative, in a compact coordinate neighborhood. Splitting
the integral into \(U\) and its complement and using (DI.7) proves
(DI.5). The same argument proves that \(A_Nv\) is smooth. It applies
to every finite list of words simultaneously and therefore to each
of the literal \(\mathrm{Sobolev}_2\) seminorms. Conjugation by the
stored frame carries (DI.8) to \(\mathrm{lieAt}\) without changing
the vectors. Thus neither a different smooth globalization nor a
different Hilbert model has been introduced.
% END DI2

% BEGIN DI3
\paragraph{DI3: absolute irreducibility and the finite interpolation argument.}
\label{claim:DI3}
For each original discrete-series model, its smooth \(K\)-finite
vectors are an irreducible module for the algebra generated by the
derived real Lie operators, and the commuting endomorphisms are
exactly the complex scalars. For finitely many such modules, a nonzero
subspace of their algebraic tensor product stable under the derived
operators in each factor is the whole tensor product.

\emph{Proof.}
In positive rank, reorder the signs of the stored Hermitian frame to
the usual \((p,q)\) ordering and pull the same unitary representation
back along the source's determinant-square-root cover. Explicitly this
is the closed subgroup
\(\{(g,z)\in U(p,q)\times\mathbb C^\times:z^2=\det g\}\).
It is locally compact, has its Borel structure and a Haar measure,
and has the two coordinate homomorphisms required by
\(\mathrm{GrpData}\). Its exponential on the real Lie algebra is
\(X\mapsto(\exp X,\exp(\operatorname{tr}X/2))\).
The identity \(\det\exp X=\exp\operatorname{tr}X\), obtained
by differentiating the determinant along \(\exp(tX)\), verifies
the defining equation; commuting scalar multiples give the
one-parameter group law. The other cover fields are coordinate
identities. Every nonzero complex determinant has a square root,
so the first projection is onto. Surjectivity
of the cover preserves irreducibility. Its compact preimage has the
same finite-orbit vectors, and its Lie algebra maps isomorphically to
the original one. The retained \(\mathrm{HarishChandraModules}\)
predicate gives an injective Harish--Chandra module map whose range
is exactly those vectors, with \(\mathrm{IrrGK}\) and
\(\mathrm{AdmGK}\). Its derivative clause identifies the action
with the original derived action. Real Lie matrices complexly span
the full Lie algebra: decompose a matrix into its fixed and
anti-fixed parts for the conjugate-linear real-form involution.
Consequently \(\mathrm{IrrGK}\) gives the asserted irreducibility.

Here \(\mathrm{AdmGK}\) states finite-dimensionality of each compact
highest-weight space. One such space is nonzero. To see this directly,
take a nonzero weight vector in a finite-dimensional compact-stable
space and maximize \(-\sum_a a\mu_a\), separately within the two
ordered sign blocks, among its finitely many weights. A compact
raising operator either kills the vector or increases that functional;
the latter is impossible. This produces a nonzero highest vector.
The weights have rational coordinates: the lifted diagonal circles
have period dividing twice the original period, so their characters
have half-integral exponents. Thus this is one of the precise
\(\mathbb Q\)-indexed highest-weight spaces in \(\mathrm{AdmGK}\).
An endomorphism commuting with all derived operators preserves its
finite-dimensional highest-weight space. Over \(\mathbb C\) it has
an eigenvector there. The kernel of \(T-c\) is then a nonzero invariant
submodule, hence the whole module. Thus \(T=c\). In rank zero the
group is trivial, and irreducibility says that any nonzero line is the
whole Hilbert space; the same conclusions hold directly. This also
checks the empty matrix case.

We give the finite interpolation argument rather than invoke a density
theorem. Let \(A\) be the unital algebra generated by the operators
on an irreducible module \(M\), and suppose
\(\operatorname{End}_A(M)=\mathbb C\). For linearly independent
\(v_1,\ldots,v_r\) and arbitrary \(w_1,\ldots,w_r\), there is
\(a\in A\) with \(av_i=w_i\). For \(r=1\) this is irreducibility.
Assume it for \(r-1\), and put
\(I=\{a\in A:av_1=\cdots=av_{r-1}=0\}\).
The subspace \(Iv_r\) is invariant. If it were zero, the rule
\[
       (av_1,\ldots,av_{r-1})\longmapsto av_r
 \tag{DI.9}
\]
would be a well-defined \(A\)-linear map \(M^{r-1}\to M\), since
the induction hypothesis makes its domain all of \(M^{r-1}\).
Each component is scalar, giving \(v_r=\sum_{i<r}c_iv_i\), a
contradiction. Thus \(Iv_r=M\); choose the first \(r-1\) images
and then correct the last using \(I\). Prescribing a linear map on
any finite-dimensional span is now possible by choosing a basis.

Let \(N\ne0\) be a factorwise-invariant tensor subspace, and choose
\(0\ne z\in N\). In each factor only finitely many coefficient
vectors occur in an expression for \(z\). The preceding interpolation
therefore permits any desired endomorphism of those finite spans to
be applied by an element of the corresponding factor algebra. Choose
linear functionals in the factors whose tensor evaluation on \(z\)
is nonzero; finite basis coordinates supply them. Applying their
rank-one maps sends \(z\) to a nonzero multiple of any prescribed
pure tensor. Hence \(N\) contains all pure tensors and is the whole
algebraic tensor product. The argument also covers an empty set of
factors, whose tensor product is \(\mathbb C\).
% END DI3

% BEGIN DI4
\paragraph{DI4: one nonzero vector separates an equivariant multiplicity map.}
\label{claim:DI4}
Let \(U\) be a linear map from the algebraic tensor product of the
original smooth models to functions on the target adelic group.
Suppose it is equivariant in every archimedean factor and its point
evaluations are separately Sobolev-continuous, in the literal
\(\mathrm{ContAt}\) sense. If \(\ker U\) contains a nonzero tensor,
then \(U=0\).

\emph{Proof.}
The kernel is invariant under each real group. It is also stable under
the corresponding derived operators. For a smooth tensor, differentiate
its orbit in one factor. Its difference quotients converge in every
fixed finite Sobolev seminorm, by applying the fundamental theorem of
calculus to one further orbit derivative. The separately continuous
point evaluations of \(U\) allow that limit to pass through \(U\).
The equivariance equation therefore differentiates, proving stability
of the kernel.

The same reasoning passes compact integrals through each point
evaluation. Thus the averages \(A_N\) of DI2, applied in one tensor
factor, preserve \(\ker U\). If \(z\ne0\) lies there, apply these
averages in all factors. Their Hilbert-tensor norms tend to \(z\),
because \(z\) is a finite sum of pure tensors and the averages are
contractions converging strongly on every coefficient vector. The
algebraic tensor product of Hilbert subspaces injects into the Hilbert
tensor product: on the finitely many coefficient spans this is the
ordinary positive definite tensor Gram matrix. Hence some averaged
tensor is nonzero. It is in the tensor product of the \(K\)-finite
subspaces and in \(\ker U\).

DI3 now puts that whole \(K\)-finite tensor product in the kernel.
To reach arbitrary smooth tensors, hold all but one factor fixed and
use DI2 and its Sobolev convergence. Repeat over the finite factor
set. Separate continuity suffices; no unproved joint continuity or
completed tensor-product theorem is used. This proves \(U=0\).

In particular the conclusion applies to
\[
             U=T(\,\cdot\,)(\,\cdot\,\kappa)-cT
 \tag{DI.10}
\]
for an original multiplicity map \(T\) and finite \(\kappa\).
The finite and infinite adelic components commute, so (DI.10) is
archimedean equivariant. At a point \(g\), its Sobolev functional
is the difference of the ones for \(T\) at \(g\kappa\) and \(g\);
concatenating their finite derivative lists gives the required bound.
This specialized statement is sufficient below. It does not assert
that an arbitrary separately continuous map is separated by an
arbitrary nonzero vector.
% END DI4

% BEGIN DI5
\paragraph{DI5: actual adelic lifts of the common residual block.}
\label{claim:DI5}
For each chain space \(V=Y_j\) and for the terminal space \(X'\),
the original frame and level data lift every
\(g_0\in G_0(P_g,r)\) to an actual finite adelic unitary element
\(\kappa\) with \(\mathrm{IsLevelLift}(V,\beta,L,g_0,\kappa)\).

\emph{Proof.}
At a selected split prime \(\mathfrak P\in P_g\), form the literal
block diagonal matrix \(\operatorname{diag}(g_{0,\mathfrak P},1)\)
of size \(\dim V\). Lift its entries to the completed local ring.
Its determinant has nonzero residue, hence is a unit; the adjugate
formula gives its integral inverse. Write this lift in the stored
\(\beta\)-frame. At the conjugate prime put the inverse Hermitian
transpose prescribed by the original Hermitian pairing. These two
matrices together are unitary. The self-duality of the framed lattice
shows that the conjugate matrix and its inverse preserve the same
completed lattice: the dual of an integral lattice automorphism acts
integrally on the dual lattice, which here is the original lattice.
Set every other finite component and every infinite component to one.
This is an element of the restricted adelic product, not a rational
approximation to one.

The original level conditions say that no prime of \(P_g\) is
conjugate to another prime of \(P_g\cup Q\), and that \(P_g\cap Q\)
is empty. Thus this construction is the identity at all other selected
residues and at every \(Q\)-residue. Taking its finite product over
\(P_g\) gives the prescribed \(g_0\) simultaneously. The block
diagonal matrix is upper block triangular for the partition
\(r,t,\ldots,t\), its top \(r\)-block is exactly \(g_0\), and it
satisfies the two-block condition at \(Q\). This checks all three
fields of \(\mathrm{IsLevelLift}\): finite support, parahoric
membership, and the exact residual block.

The source's \(\mathrm{IsSDFrame}\) supplies the full self-dual
lattice and the stored frame used here; its exponent \(1^a\) equals
one. The same proof applies in rank zero or when the complementary
block is empty. In particular it applies at \(j=m\), the terminal
rank \(r+mt\), as well as every inner rank \(r+jt\). A residue
lift chosen for a particular \(g_0\) is used only in the proof of a
universal identity and does not replace any of the selected readers.
% END DI5

% BEGIN DI6
\paragraph{DI6: the two exact finite datum transformation laws.}
\label{claim:DI6}
Write \(a=r+jt\), \(b=r+(j+1)t\), and let \(G,H\) be the
reductions of matching finite level lifts on these two spaces. At
\(P_g\) they are upper block triangular for the prescribed blocks,
with the same upper \(r\)-block. Let
\[
 A_j(G)=\prod_{c=1}^{j}\chi_c(\det G_{cc}),\qquad
 A_{j+1}(H)=\prod_{c=1}^{j+1}\chi_c(\det H_{cc}).
\]
The original type datum transforms under \(y\mapsto H^{-1}yG\) as
\[
 \Omega_G(H^{-1}yG)=
  A_j(G)A_{j+1}(H)^{-1}
  \chi_{j+1}(\det H)\chi_{j+1}(\det G)^{-1}\Omega_G(y).
 \tag{DI.11}
\]
At a twist prime, put \(G_{\rm lo},H_{\rm lo}\) for the lower
diagonal blocks of the partition \(r+(\dim V-r)\). The actual
twist datum, whose exponent in \(\mathrm{IsStepVecOf}\) is \(+1\),
satisfies
\[
 \Omega_Q(H^{-1}yG)=
 \xi(\det G)\xi(\det G_{\rm lo})^{-1}
 \xi(\det H)^{-1}\xi(\det H_{\rm lo})\Omega_Q(y).
 \tag{DI.12}
\]
Both equations include the zero-support cases.

\emph{Proof.}
For the type datum its support is precisely: all rows below \(a\)
vanish, the upper square block \(Y\) is invertible and upper block
triangular, and its upper \(r\)-block equals one. Multiplication
gives the new square block
\[
                       Y'=H_{\le a}^{-1}YG .
 \tag{DI.13}
\]
The common upper \(r\)-blocks cancel, so its top block is again one.
All support conditions are preserved in both directions by the same
calculation with inverse matrices. Each diagonal-block determinant
of \(Y'\) is the corresponding determinant of \(Y\), multiplied
by that of \(G\) and divided by that of \(H\). The definition is
\(\Omega_G(Y)=A_j(Y)\chi_{j+1}(\det Y)^{-1}\).
Use
\(\det H=\det H_{\le a}\det H_{j+1,j+1}\) and collect the
\(j+1\) factors. This is (DI.11), with no suppressed determinant.

For the twist datum, the condition is that the lower rows of the
first \(r\) columns vanish and the top \(r\)-square determinant
is nonzero. Its top block becomes \(H_{00}^{-1}Y_{00}G_{00}\).
The support is again preserved in both directions. Applying \(\xi\)
to its determinant and using
\(\det G=\det G_{00}\det G_{\rm lo}\) and the analogous formula
for \(H\) gives (DI.12). If a support condition fails, it fails
after transformation as well and both sides are zero. Products over
the actual finite prime sets give the source's complete
\(\mathrm{omegaG}\) and \(\mathrm{omegaQ}\) laws. No character
of an arbitrarily chosen finite group has been substituted.
% END DI6

% BEGIN DI7
\paragraph{DI7: the actual step vacuum has the flag-character ratio.}
\label{claim:DI7}
Let \(\kappa,\kappa'\) be matching finite lifts at adjacent chain
spaces, including the terminal pair \((Y_{m-1},X')\). Put
\[
 a_0=\mathrm{levChi}(\mathrm{flagL}(L_j))(\kappa),\qquad
 a_1=\mathrm{levChi}(\mathrm{flagL}(L_{j+1}))(\kappa').
\]
For the original stored \(v=vY(t,j)\), or \(v=vX(t)\) at the
last step, one has
\[
                 \omega(\kappa,\kappa')v=(a_0/a_1)v.
 \tag{DI.14}
\]

\emph{Proof.}
We first justify the finite-model transport on the given functions.
The actual step half-function has the required Schwartz--Bruhat
guard. Indeed its support is contained in the chosen finite half,
and the integrality clauses bound it in a fractional lattice.
Choose one positive integer that clears the finitely many
denominators in the original frames and weighted lattices, including
the fractional-ideal factor and its inverse. Enlarge that integer
by the finitely many selected residue primes. Translation by the
resulting small lattice preserves the selected half, every
integrality test, and every residue matrix occurring in
\(\mathrm{omegaG}\) and \(\mathrm{omegaQ}\). It therefore
preserves both their characters and their zero branches. Enlarging
the integer once more places the whole support in its inverse
lattice. These are precisely the support and periodicity clauses
of \(\mathrm{IsSBHalf}\). This is the same finite denominator
argument as FIA4, applied here to the step function with its
displayed positive twist exponent and its own chosen half.
Let \(\Lambda=\mathrm{latO}(S,T,\Lambda_S,\Lambda_T,\mathrm{True})\)
be the original weighted Hom lattice, including its fractional-ideal
factor. A finite half-function with the \(\mathrm{IsSBHalf}\) guard
is supported in \(k^{-1}\Lambda_e\) and periodic by
\(k\Lambda_e\) for some \(k>0\). It is a finite combination of
coset indicators. On each such finite quotient, the exact character
orthogonality formula
\[
  1_{k\Lambda_e}(x)=\frac1{[\Lambda_e:k\Lambda_e]}
       \sum_{\alpha}\alpha(x)1_{\Lambda_e}(x)
 \tag{DI.15}
\]
expresses the indicator as modulations of \(1_{\Lambda_e}\),
followed by a translation for a nonzero coset. Self-duality realizes
these characters by the opposite finite half: the annihilator of
\(\Lambda_e\) is \(\Lambda_{\bar e}\), and scaling by \(k\)
gives the corresponding perfect finite pairing. This is the finite
pairing calculation of WA12 and QD19, with no coefficient bound needed
here. Consequently every place model with the original translation,
modulation and vacuum axioms sends every guarded half-function to a
finite linear combination of finite Heisenberg translates of its
vacuum. These are theta functions, by the original Heisenberg
quasi-periodicity and smoothness calculation in JIO13.

If \(A=\operatorname{actO}(\kappa,\kappa')\), covariance and
stabilization of \(\Lambda\) show that
\[
             T'=c^{-1}\omega(\kappa,\kappa')\,T\,A^*
 \tag{DI.16}
\]
is another place model with the same vacuum whenever
\(\omega(\kappa,\kappa')f_0=cf_0\), \(c\ne0\).
Here \(A^*\Psi=\Psi\circ A\). To check (DI.16), substitute it
in the two half-translation laws; symplectic invariance moves the
modulation pairing through \(A\), and Weil--Heisenberg covariance
moves the translation. The vacuum indicator is fixed by \(A^*\),
so the normalization is exactly one. The preceding paragraph verifies
the theta-function guard every time covariance is applied. The
universal part of the original \(\mathrm{IsThetaAt}\), together
with \(\Phi\circ A^{-1}=c'\Phi\), then gives
\(\omega(\kappa,\kappa')v=cc'v\). This proves transport of the
stored vector itself, even if that vector is zero.

The self-dual frame vacuum law BTH15 applies to the original full
Weil action. At \(P_g\) its factor is
\(\chi_{j+1}(\det G)\chi_{j+1}(\det H)^{-1}\); at \(Q\) it is
\(\xi(\det G)^{-1}\xi(\det H)\).
Here the exact residue restrictions come from the original
\(\mathrm{LevCh.Is}\): both consecutive splitting characters have
inverse \(\chi_{j+1}\) residue at \(P_g\) and \(\xi_t\) residue
at \(Q\), and the remaining \(\eta\)-factors are unramified on
finite units. For completeness, the character restriction to the
base field is the original quadratic character to the rank power;
adjacent ranks differ by the original odd \(t\). Dividing the two
restrictions gives the required \(\eta\) restriction, since that
quadratic character has square one. Its stored infinity law and
local-unit triviality give the exact \(\mathrm{IsEtaTilde}\)
normalization used in BTH15; finite-support integral units are dense
in their compact product. Thus this use retains the original
\(\eta\), not a new splitting character.

Multiplying these vacuum factors by (DI.11)--(DI.12) cancels every
full determinant. What remains is
\[
 \frac{\displaystyle
   \prod_{P_g}A_j(G)\prod_Q\xi(\det G_{\rm lo})^{-1}}
 {\displaystyle
   \prod_{P_g}A_{j+1}(H)\prod_Q\xi(\det H_{\rm lo})^{-1}}
                          =a_0/a_1 .
 \tag{DI.17}
\]
The frames intertwine \(\operatorname{actO}^{-1}\) with
\(y\mapsto H^{-1}yG\), preserve integrality and the selected half,
and induce exactly these residue matrices. Outside support the
finite function is zero on both sides. This verifies the whole
\(\Phi\) used in \(\mathrm{IsStepVecOf}\).
The proof uses only the adjacent-rank equalities and applies without
change to \(j=m-1\), where \(\dim X'=r+mt\). Thus (DI.14) includes
the last step, not only an inner step with \(j+1<m\).
% END DI7

% BEGIN DI8
\paragraph{DI8: the same character on every original polynomial vector.}
\label{claim:DI8}
If a finite pair has \(\omega(\kappa,\kappa')v=cv\), then, for
every polynomial \(P\) in the original Fock variables,
\[
           \omega(\kappa,\kappa')\mathrm{pt}(P)
                      =c\,\mathrm{pt}(P).
 \tag{DI.18}
\]
For a source level element \(\kappa\) at either \(P_b\), DI5
provides a matching target element \(\kappa'\) at the same \(P_b\),
and here
\[
                 c=\chi_0(\kappa)\chi_1(\kappa')^{-1},
 \tag{DI.19}
\]
where \(\chi_0,\chi_1\) are the two full original level characters.

\emph{Proof.}
A finite pair is identity in every infinite matrix component. Its
symplectic action therefore fixes every archimedean Heisenberg
direction. The covariance law says it commutes with every such
Heisenberg translation. To differentiate this law, use the original
Haar-norm identity: \(\omega\) is an isometry on the theta
\(L^2\) space of the compact rational additive quotient. The
real Heisenberg difference quotients converge there by CW12's
Banach smoothness calculation. Applying the isometry preserves their
limits. Both resulting derivatives are continuous theta functions;
equality in \(L^2\) is pointwise equality because a nonzero
continuous function has positive norm on some open set of positive
Haar measure. Thus \(\omega\) commutes with each actual scaled
creator. Induct on the ordered creation word defining a monomial,
then use linearity. This keeps all creator constants and the original
vacuum, and proves (DI.18).

For (DI.19), let \(g_0\) be the upper residual block of the given
\(\kappa\). DI5 chooses \(\kappa'\) with the same block. Since
\(g_0\in P_b.H\), this target lift lies in the desired full level.
Each full character is \(P_b.\chi(g_0)\) times its flag character.
The common nonzero factor cancels in their ratio. Apply DI7 to obtain
(DI.19), including the terminal pair and both values of \(b\).
% END DI8

% BEGIN DI9
\paragraph{DI9: exact translation of the stored theta integral.}
\label{claim:DI9}
Consider an original adjacent pair of the chain, with target map
\(y\), output \(x=\mathrm{io}(y)\), and real intertwiner
\(\mathrm{It}\). If \(y\) has target level character \(\chi_1\)
and the polynomial kernel obeys (DI.18), then
\[
 x(\mathrm{It}(P,w))(g\kappa)
    =c\chi_1(\kappa')x(\mathrm{It}(P,w))(g)
 \tag{DI.20}
\]
for every polynomial, every input smooth tensor, and every \(g\).
The statement uses the source measure and fundamental domain stored
in that particular step.

\emph{Proof.}
The literal \(\mathrm{IsIo.eq}\) is
\[
 x(\mathrm{It}(P,w))(g)
   =\int_{\mathcal D}\theta_\omega(g,h,\mathrm{pt}(P))\,
                                      y(w)(h)\,d\nu(h).
 \tag{DI.21}
\]
There is no conjugation of \(y\) in this equation. The representation
law and (DI.18), with
\((g\kappa,h)=(g,h\kappa'^{-1})(\kappa,\kappa')\), give
\[
 \theta_\omega(g\kappa,h,\mathrm{pt}(P))
      =c\,\theta_\omega(g,h\kappa'^{-1},\mathrm{pt}(P)).
\]
The target level law gives
\(y(w)(h)=\chi_1(\kappa')y(w)(h\kappa'^{-1})\).
Both factors in the resulting integrand are left invariant under the
same rational target group, by automorphy of the original Weil datum
and of \(y\). RQC4 therefore identifies its integral over
\(\mathcal D\) with that of its right translate by
\(\kappa'^{-1}\). The Haar measure is right invariant, and translating
the fundamental domain introduces no volume factor. This proves
(DI.20). Linearity extends the formula from pure input tensors to
all algebraic input tensors.

For inner steps the measure is exactly \(R.\nu_Y(t,j+1)\), with
the final inner value read from \(\mathrm{models}_{\nu Y'}\) when
\(j+1=m-1\); the other values are read from
\(\mathrm{models}_{\nu Y}\). For the last step it is
\(R.\nu_X(t)\), with \(R.\mathrm{fdX}(t)\).
These are the original Haar/fundamental-domain clauses, so no
identification of two independently normalized measures is needed.
% END DI9

% BEGIN DI10
\paragraph{DI10: the literal inner and terminal level-preservation predicates.}
\label{claim:DI10}
The data of DI1 satisfy both
\[
 \mathrm{StepLev}(R.\mathrm{toRealDat},F.\mathrm{lv2}),\qquad
 \mathrm{TopLev}(R.\mathrm{toRealDat},F.\mathrm{lv2}).
 \tag{DI.22}
\]

\emph{Proof.}
Fix \(b,t,j\), with \(j+1<m\), and an arbitrary
\(y\in R.CC(t,j+1)\) belonging to the full \(P_b\)-level
multiplicity subspace. Its image \(x=R.\mathrm{ioY}(t,j)y\)
already belongs to the original principal multiplicity space, by its
codomain. Thus automorphy, cuspidality, archimedean equivariance and
separate Sobolev continuity are already properties of \(x\).
Fix any \(\kappa\) in the desired source level. DI8 gives
\(\kappa'\) and \(c=\chi_0(\kappa)/\chi_1(\kappa')\).
All level characters take unit values and in particular are nonzero.
DI9 says that the defect
\[
                     U=x(\,\cdot\,)(\,\cdot\,\kappa)
                                      -\chi_0(\kappa)x
\]
annihilates every \(\mathrm{ItY}(j)(P,w)\). The original
\(\mathrm{IsItW.ne}\) ensures at least one such value is nonzero.
Indeed if all values on pure input tensors were zero, multilinearity
and their span would make the entire stored map zero. DI4 applies
to this defect and gives \(U=0\). This is the full level equation,
for all output smooth tensors and all \(g\), and proves
\(\mathrm{StepLev}\).

For \(\mathrm{TopLev}\), start with any element of the literal
range
\[
        \operatorname{range}\bigl(R.cX.\mathrm{toMap}(R.prX)
                                      \circ R.HX(t).e\bigr)
 \tag{DI.23}
\]
which has the full terminal \(P_b\)-level. Its stored
\(R.\mathrm{ioX}(t)\) image lies in \(R.CC(t,m-1)\).
Repeat the preceding proof with \(\omega_X,v_X,\mathrm{ItX},
\nu_X,\mathrm{fdX}\), and DI7's terminal character identity.
The terminal \(\mathrm{IsItW}\) again supplies nonzeroness on the
original target tensor space. DI4 now proves the whole source level
law. No enlargement of (DI.23) to an arbitrary source representation
is made, and no new real intertwiner is selected.
% END DI10

% BEGIN DI11
\paragraph{DI11: the actual finite row idempotent and its adjoint.}
\label{claim:DI11}
The original \(e_0\) preserves each full \(P_b\)-level
multiplicity space, maps it into its good part, and is self-adjoint
for the actual quotient integrals of the row pairs in (DI.3a),
under that retained hypothesis.

\emph{Proof.}
On the principal flag space the action of a residual element \(g\)
is the literal operator
\[
             A_gf(h)=a_g^{-1}f(h\kappa_g),
 \tag{DI.24}
\]
where \(\kappa_g\) is any lift from DI5 and \(a_g\) is its flag
character. Changing lifts changes the right translate by exactly the
principal flag character, so (DI.24) is independent of the lift.
Multiplying lifts proves \(A_gA_h=A_{gh}\), with the source's
right-translation convention. The stored group record gives
\[
                 e_0=\sum_{g\in G}e(g)A_{R.e_Gg},
 \qquad e*e=e,\quad e\text{ central},\quad
                     e(g^{-1})=\overline{e(g)}.
 \tag{DI.25}
\]
These last identities are the finite central character idempotent
of the actual selected block. Its coefficients and the group
identification are those in \(\mathrm{Pins}.\mathrm{keep}.\mathrm{grp}\).
Finite convolution of (DI.25) proves \(e_0^2=e_0\) and commutation
with every \(A_g\). A function in the full \(P_b\)-level is
character-eigen under the \(A_p\), \(p\in P_b.H\), because its
full level character is \(P_b.\chi(p)a_p\). Conversely these
eigenrelations and its principal flag level imply its full level
law by (DI.24). Hence \(e_0\) preserves that full level, and its
image is fixed by \(e_0\), which is exactly the good-part condition.
Finite translates and linear combinations preserve the other four
multiplicity-map clauses; the Sobolev lists concatenate.

All flag-character values in (DI.24) have modulus one. This follows
from their prescribed residue characters, or from finite order in
the finite residual group. The inverse lift has character
\(a_{g^{-1}}=a_g^{-1}\). RQC4, with the original fundamental
domain, therefore gives for integrable products
\[
 \int_{\mathcal D}\overline{A_gf}\,h\,d\mu
       =\int_{\mathcal D}\overline f\,A_{g^{-1}}h\,d\mu.
 \tag{DI.26}
\]
For example, in the left integral substitute \(u=h_0\kappa_g\):
the scalar \(\overline{a_g^{-1}}=a_g\) is exactly the inverse
operator's scalar. Reindex \(g\mapsto g^{-1}\) in the finite sum
and use the last identity in (DI.25). This proves
\[
         \int\overline{e_0f}\,h\,d\mu
                    =\int\overline f\,e_0h\,d\mu.
 \tag{DI.27}
\]
For the row pairs used below, (DI.3a) supplies each left-hand
integrable product in (DI.26). Its change of variables proves
integrability of the right-hand product as well. Finite summation
then gives the integrability needed on both sides of (DI.27),
on exactly \(R.\mathrm{fd0}\) with exactly \(\mu_b\).

The analytic boundary is substantive. An arbitrary original
\(\mathrm{IsAutFunSm}\) supplies one exponent for its derivative
bounds on each finite coset, in \(\mathrm{normInf}\). The
source predicate \(\mathrm{CuspRapidDecay}\) instead requires
\(\mathrm{IsUnifModerate}\), a bound in the global adelic height.
The forward direction of \texttt{ExistsSigmaList.growth\_iff}
uses a polynomial unitary reduction theorem to pass between them.
That conversion has not been reconstructed in this chapter. If it
is supplied, applying \(\mathrm{CuspRapidDecay}\) with exponent
zero gives boundedness; GV10's finite quotient volume then implies
(DI.3a). This explains the source route but does not discharge the
retained analytic supplier by invoking that route's theorem names.
% END DI11

% BEGIN DI12
\paragraph{DI12: the literal row-projector metric, including its choice definition.}
\label{claim:DI12}
The data of DI1 satisfy
\(\mathrm{RowProj}(R.\mathrm{toRealDat},F.\mathrm{lv2})\).

\emph{Proof.}
Fix \(b,t\), let \(w\) be in the full level multiplicity space,
and let \(c\) lie in its good part. By definition
\(\mathrm{postOn}(e_0)c=c\). DI11 gives
\(\mathrm{postOn}(e_0)w\) in that same good part. For arbitrary
smooth tuples \(v,v'\), apply (DI.27) to \(c(v),w(v')\). Since
\(e_0c(v)=c(v)\), it gives
\[
 \int_{\mathrm{fd0}}\overline{c(v)}w(v')\,d\mu_b
  =\int_{\mathrm{fd0}}\overline{c(v)}
                      (\mathrm{postOn}(e_0)w)(v')\,d\mu_b.
 \tag{DI.28}
\]
Both products are integrable by DI11.

The literal definition of \(\mathrm{ipM}(c,w)\) chooses a scalar
\(z\) satisfying, for \emph{all} \(v,v'\),
\[
 \int_{\mathrm{fd0}}\overline{c(v)}w(v')\,d\mu_b
          =z\prod_\tau\langle v_\tau,v'_\tau\rangle;
 \tag{DI.29}
\]
if no such scalar exists, its value is zero. Equations (DI.28) and
the integrability assertions prove that exactly the same scalars
satisfy (DI.29) for \((c,w)\) and \((c,\mathrm{postOn}(e_0)w)\).
If none exist, both definitions are zero. If they exist, the scalar
is unique: choose a nonzero smooth vector in every local factor,
take \(v'=v\), and divide by the positive product of their squared
norms. Such vectors exist by the nontrivial irreducible models and
the nontrivial smooth Harish--Chandra module in DI3, with its
separate rank-zero case; the empty factor product is one.
Thus the two choices give the same number, proving exactly
\[
 \mathrm{ipM}_{\mu_b}(c,w)
       =\mathrm{ipM}_{\mu_b}(c,\mathrm{postOn}(e_0)w).
 \tag{DI.30}
\]
This proof does not assume the existence of a Schur scalar for every
pair. It proves the original definition's equality even in its
zero fallback branch. Its positive scalar measure \(n_b^{-1}\nu_0\)
has remained fixed on both sides throughout.
% END DI12

% BEGIN DI13
\paragraph{DI13: the entire original composite intertwiner.}
\label{claim:DI13}
For the fixed readers in DI1, the following two equalities hold for
every original tuple \(t\) and every row-coordinate vector \(x\):
\[
\begin{split}
 (Y.I_{\rm sh}.j_E(D.\mathrm{what}_2.\iota(Y.\mathrm{tup}(t))
                                      (Y.j_A(t)x)))_{\rm val}
 &=\mathrm{lift}\bigl(R.HE.e\bigl(R.\mathrm{Th}(t)
                  (R.j_X(t)\mathrm{rowForm}_1(t,x))\bigr)\bigr),\\
 (Y.I_0.j_E(X_0.s.\iota(Y.\mathrm{tup}(t))
                                      (Y.j_{A0}(t)x)))_{\rm val}
 &=\mathrm{lift}\bigl(R.HE.e\bigl(R.\mathrm{Th}(t)
                  (R.j_X(t)\mathrm{rowForm}_0(t,x))\bigr)\bigr).
\end{split}
 \tag{DI.31}
\]
Here \(\mathrm{lift}\) is the original linear equivalence between
multilinear maps and linear maps on the algebraic tensor product.
\(\mathrm{rowForm}_b\) uses exactly \(F.P_b,F.\psi_b\).

\emph{Proof.}
We first record the finite Hilbert-space argument underlying the
whole chain. Suppose \(j:H\simeq C\subset N\) identifies its
inner product with the given pairing \(B\) on \(C\), and a map
\(f\) is specified by
\[
             \langle z,f(y)\rangle=B(jz,g(y)).
\]
If \(p(y)\in C\) and \(B(c,g(y))=B(c,p(y))\) for all \(c\in C\),
then \(jf(y)=p(y)\): test against every \(z\), and use
nondegeneracy of the Hilbert inner product. In particular, when
\(g(y)\in C\), its orthogonal projection is itself. This is the
exact argument for the source's \(\mathrm{HilbTie}\) and
\(\mathrm{ProjTie}\), without replacing either pairing.

Start at the original terminal reading \(Y.\mathrm{topRead}\),
or \(Y.\mathrm{topRead0}\) for the base side. It identifies the
input with the stored vector in (DI.23). DI10 puts its terminal
image in the required full level. The original \(\mathrm{ioTop}\)
projection identity therefore equals that image itself. Descend
through \(j=m-2,\ldots,1\). At each step DI10 again supplies full
level membership, and the same projection argument identifies the
abstract row \(\mathrm{io}\) with the concrete \(R.\mathrm{ioY}\).
At the last descent to level zero, use DI12: the actual projected
value is \(\mathrm{postOn}(e_0)\) of the concrete descent, for the
pairing \(\mathrm{ipM}_{\mu_b}\) specified by \(Y.\mathrm{tied}\).
The induction is valid for \(m\ge2\), which the original ranks
provide. If written with the more general \(m=1\) convention, the
single terminal-to-zero projection is precisely the same last
argument; no empty interior product is assigned an extra factor.

Unfolding \(\mathrm{toChainReal}\) gives, at every \(j<m\), the
same \(R.CC(t,j)\) and the same \(\mathrm{ioY}(t,j)\) value.
Its terminal map is the same \(\mathrm{ioX}(t)\) on the range
(DI.23). Thus its descending composition is literally
\(R.\mathrm{thDn}(t)\), not merely an isomorphic composition.
The full chain's stored \(\mathrm{Th}\) equation is
\[
              R.HE.e(R.\mathrm{Th}(t)y)
                         =e_0(R.\mathrm{thDn}(t)y).
 \tag{DI.32}
\]
It is active by DI1. This identifies the bottom of the induction
with the right side of (DI.31).

Finally every \(Y.\mathrm{tup}(t)\) is a used tuple. The original
\(\mathrm{mu\_used}\) field says either \(t=\mu_0\), where
\(Y.\mathrm{tup0}\) applies, or \(t=R.\mu(w)\), where
\(Y.\mathrm{muAt}\) and the original prime equivalence apply.
The left inverse \(Y.\mathrm{tupT\_tup}\) identifies the chain
index with \(t\). For the sharp side an induction on the number of
steps identifies \((\mathrm{RowSide.sharp}\,D).\mathrm{downTo}\)
with \(D.\mathrm{what}_2.\mathrm{downTo}\), so its \(\iota\) is
the one displayed in (DI.31). This proves both identities on the
actual readers and their two original metrics.
% END DI13

% BEGIN DI14
\paragraph{DI14: the exact quotient-coordinate projection inside theta.}
\label{claim:DI14}
Let \(X=R.\mathrm{toCx3}(F)\), and let \(b=0\) or \(1\).
For the actual level row embedding \(e_V\), actual quotient
coordinates \(F.cr.\mathrm{toR}\), and actual matrix \(P_q\),
one has
\[
 \Theta_t j_R(e_Vx)
     =\Theta_t j_R\bigl(F.cr.\mathrm{toR}(P_qx)\bigr).
 \tag{DI.33}
\]
The statement retains the original positive theta operators and
rational row tests in \(\mathrm{Holds}.\mathrm{prim}\).

\emph{Proof.}
We check exactly which earlier arguments supply the required spaces.
The \(\mathrm{models}\) clauses identify the original coefficient
form \(j_R\) and its Hermitian metric; the conjugate \(\mathrm{Tw}\)
convention applies the same involution to both entries and gives the
literal \(\mathrm{JIsoR}\) equality. The primitive
\(\mathrm{rowTest}\) and that coefficient form give the actual
rational adjoint preimage by TD1--TD3. Its row Hecke-adjoint and
group-equivariance clauses preserve the kernel and its orthogonal
complement by TD4--TD5. Apply TD6--TD7 to the exact finite level
idempotent, using \(\mathrm{Ok3}\) for its selected character block
and \(\mathrm{levE}\) for its actual target character. They give
\[
\begin{split}
 \operatorname{span}_{\mathbb C}j_R(F.cr.\mathrm{toR}(L^{d_Q}))
       &=\operatorname{im}\Theta_t^*\cap V_{0,b},\\
 \Theta_t V_{0,b}&\subset E_b,\qquad
                       \Theta_t^*E_b\subset V_{0,b}.
\end{split}
 \tag{DI.34}
\]
These are precisely \(\mathrm{RatLev,DnLev,UpLev}\) on the
same \(X\). The \(\mathrm{sumE,posScal,thmB}\) fields give the
positive operator equations of TD8. TD9, with the same
\(\mathrm{rowTest}\) for all tuples, identifies the actual
\(\mathrm{rowMat}\), hence the \(\iota\) matrix in the stored
\(S_b= A_3.\mathrm{indefDown}(s_b\mathrm{Of}(R,F))\).
No lift-test, \(\mathrm{ratG}\), or \(\mathrm{MLat}\) existence
is inferred by this matching.

Here is the remaining projection argument. Put
\(q=F.cr.\mathrm{toR}(P_qx)\) and \(d=e_Vx-q\).
Both terms lie in the actual rational level space. Its exact
coefficient metric makes \(d\) orthogonal to every rational
quotient coordinate. To verify this last assertion directly, write
\(V\) for the level embedding matrix, \(Z\) for the rational-cut
matrix, and \(H\) for the ambient Gram matrix. The original
definitions give
\[
 V^*HV=n_bH_V,\qquad H_Z=Z^*H_VZ,
       \qquad P_q=H_Z^{-1}Z^*H_V.
 \tag{DI.35}
\]
The Gram matrix \(H_Z\) is invertible because the original block
form is positive and its chosen coordinate basis is independent;
this also covers the zero-dimensional matrix. Therefore
\[
       (VZ)^*HV\,ZP_q=(VZ)^*HV.
 \tag{DI.36}
\]
This proves the claimed pairing equality, with its exact \(n_b\)
factor. At \(b=1\) the map \(e_V\) is the stored \(F.eV\); at
\(b=0\) it is the original basis inverse composed with \(V\).
Both are exactly the maps in (DI.35).

By (DI.34), \(\Theta_tj_Rd\in E_b\). For any \(e\in E_b\),
the adjoint vector \(\Theta_t^*e\) lies in the left-hand span in
(DI.34), and hence is orthogonal to \(j_Rd\). The adjoint identity
gives \(\langle e,\Theta_tj_Rd\rangle=0\). Taking
\(e=\Theta_tj_Rd\) proves \(\Theta_tj_Rd=0\), which is (DI.33).
This finite-dimensional argument never deduces a kernel assertion
from equal dimensions alone.
% END DI14

% BEGIN DI15
\paragraph{DI15: both literal \(\mathrm{RowSideDef3}\) records.}
\label{claim:DI15}
The following two records hold on (DI.1):
\[
\begin{split}
 &\mathrm{RowSideDef3}(D,A_3.A,\mathrm{RowSide.sharp}(D),
             F.P_1,F.\psi_1,Z_s,S_1,Y.\mathrm{tup},Z.e_{RR},Y.j_A),\\
 &\mathrm{RowSideDef3}(D,A_3.A,X_0.s,
             F.P_0,F.\psi_0,Z_s,S_0,Y.\mathrm{tup},Z.e_{RR0},Y.j_{A0}),
\end{split}
 \tag{DI.37}
\]
where \(Z_s=A_3.\mathrm{fam3.toSub2}(\mathrm{zsOf}(R,F))\).

\emph{Proof.}
The forms and ambient-lattice fields are exactly the already supplied
\(Y.\mathrm{formsA},Y.\mathrm{latA}\) and
\(Y.\mathrm{formsAtBottomLevel},Y.\mathrm{latA0}\). They are not
new isometries chosen after the comparison.

For the third field, the actual block-coordinate definitions give
\[
 \mathrm{rowForm}_b(t,rV_{0,b}(t)x)
             =\mathrm{basis}^{-1}\bigl(V_b(t)x\bigr).
 \tag{DI.38}
\]
Indeed multiply the level-coordinate matrix by \(rV_0\): the
original finite block projector is identity on the selected block
and character subspace, so its coordinate left inverse cancels.
Equivalently, multiplying their matrices gives the original
\(\mathrm{lf\_lvl\_vE}\) identity. DI14 then replaces the right
side of (DI.38), inside theta, by its \(P_q\)-projection.

By TD9's exact common-image matrix, evaluated on the basis and
extended semilinearly, this projected theta value is
\[
 \mathrm{Tw.out}\left(
     X.i_C\bigl(\sigma((S_b.\iota(t))P_{q,b}(t)x)\bigr)\right).
 \tag{DI.39}
\]
There is no new comparison isomorphism in (DI.39): uniqueness in
the original row-coordinate solve is exactly the statement
\(X.\mathrm{key\_row}\), and expansion
\(\sum_i\sigma(a_i)j(e_i)=j(a)\) is
\(X.\mathrm{iC\_cpt'}\).
DI13 identifies (DI.39) with the image of the original composite
\(\iota\) applied to \(Y.j_A(rV_0x)\), or \(Y.j_{A0}(rV_0x)\).
The stored \(Z.\mathrm{jEE}\), respectively
\(Z.\mathrm{jEE0}\), identifies that same vector with
\[
        X_b.j_{EE}\bigl(Z.e_{RR,b}((S_b.\iota(t))P_{q,b}(t)x)\bigr)
\]
after the same \(Y.I_b.j_E\) embedding. That embedding is injective,
so the two vectors are equal. This is the third field of each
record in (DI.37), for every tuple and every original level vector.
The proof treats both levels with their own positive normalization
\(n_b\); it does not identify their Hermitian forms with each other.
% END DI15

% BEGIN DI16
\paragraph{DI16: the exact finite cardinalities, ranks and distinguished place.}
\label{claim:DI16}
The original pins imply
\[
\begin{gathered}
 |G|=D.\mathrm{what}_2.\mathrm{restRowsFiniteGroupCard},\qquad
 |F.P_0|=|F.P_1|(g_0g_1)^{kf},\\
 A_3.A.r=D.R.r,\qquad A_3.A.t=D.R.t,\qquad
 A_3.A.v_0=\mathrm{InfinitePlace.mk}(\sigma).
\end{gathered}
 \tag{DI.40}
\]

\emph{Proof.}
The literal group equivalence \(R.e_G.\mathrm{trans}(c_G)\)
identifies \(G\) with the product over the two distinguished prime
fibres of \(\mathrm{GL}_r(\mathbb Z/g_i\mathbb Z)\), with exactly
\(f\) factors in each fibre. Taking cardinalities gives the first
identity by the original definition of the row group cardinality.

For the second, use the \(\mathrm{LvlTie}\) inside
\(\mathrm{Pins}.\mathrm{keep}.\mathrm{psi}\). The base group is
the product of upper unitriangular groups. After the exact
\(\mathrm{symPos}\) permutation, the sharp group has the same
free entries except the \(k\) entries in column \(r_S+k\) and
rows \(r_S,\ldots,r_S+k-1\), which must vanish. The inequality
\(r_S+k<r\) follows from the original \(m\ge2\).
The permutation is injective by checking its two index intervals,
hence bijective on the finite index set.

For a residue field of size \(q\), assigning arbitrary values to
the allowed strict-upper entries, diagonal entries one, and lower
entries zero is a bijection onto the indicated group. The determinant
is one and the entries recover the assignment. The full group has
\(q^N\) elements and the sharp subgroup \(q^{N-k}\), since the
deleted entries are exactly indexed by \(0,\ldots,k-1\).
Multiply these identities over all actual selected primes. The
residue-field equivalences in the same \(\mathrm{LvlTie}\) give
\(q=g_i\), and the two fibres each have size \(f\). The ratio
is \(\prod_{i=0}^1 g_i^{kf}=(g_0g_1)^{kf}\), proving the stated
equality of integers, without a quotient-cardinality divisibility
assumption.

Finally the original \(\mathrm{Bridge.FamTie}\) witness inside
\(\mathrm{Pins}.\mathrm{char}.\mathrm{base}.\mathrm{fam}.\mathrm{tie}\)
has its exact rank and base-place fields. Projecting those fields
gives the last three assertions. The existential witness is used
only for these equalities; it does not replace the coefficient maps
already stored in \(R\).
% END DI16

% BEGIN DI17
\paragraph{DI17: integral descent of the actual good subspace.}
\label{claim:DI17}
Let \(K\subset L\) be the original number-field extension and
\(A=R.FW(i)\). From its supplied \(\mathrm{IsKsForm}\), stability
of \(A.W_K\) under \(A.\rho\), and the original coefficient-field
values of \(\chi,\psi_0^{-1},e\), one obtains
\[
 \mathcal O_L\!\operatorname{-span}(M_K\cap W_{\rm good})
            =(\mathcal O_L\!\operatorname{-span}M_K)
                                      \cap W_{\rm good}=M_j.
 \tag{DI.41}
\]
This equality concerns the original norm-defined \(M_K\), not an
arbitrary full lattice.

\emph{Proof.}
Choose the basis supplied by \(\mathrm{IsKsForm}\): it is a
\(K\)-basis of \(W_K\) and an \(L\)-basis of \(W_L\).
For any \(K\)-linear \(\varphi:L\to K\), apply \(\varphi\)
to these coordinates to get \(T_\varphi:W_L\to W_K\).
It commutes with every \(K\)-matrix. In particular it commutes
with the group operators and with
\[
               E=\sum_g e(g)\chi(g)^{-1}\rho(g).
 \tag{DI.42}
\]
The type equations have coefficients \(\chi(p)\psi_0(p)^{-1}\)
in \(K\), so they are preserved by \(T_\varphi\). If \(z=Ew\),
then \(T_\varphi z=E T_\varphi w\). Hence \(T_\varphi\) also
preserves the image condition in \(W_{\rm good}\). This proves
\(T_\varphi(W_{\rm good})\subset W_{\rm good}\) directly for
the literal intersection of type space and image in the source.

We recall the finite algebra needed to choose integral coordinate
functionals. The finite torsion-free \(\mathcal O_K\)-module
\(\mathcal O_L\) is projective. A short local proof suffices:
at every nonzero prime of the Dedekind ring \(\mathcal O_K\),
its localization is a finite torsion-free module over a discrete
valuation ring. Successively choosing a vector of minimal valuation
and eliminating its coordinate proves that module is free. Its local
basis and dual basis extend after inverting one element outside the
prime, since there are only finitely many generators and relations.
These neighborhoods, together with the generic one, have a finite
subcover of \(\operatorname{Spec}\mathcal O_K\). Clearing the
finitely many denominators in each local dual-basis identity gives
\(f_a^{N_a}\mathrm{id}\) as a finite sum of integral rank-one
maps. The \(f_a^{N_a}\) generate the unit ideal, because no prime
contains them all. A linear combination of these identities gives
a global dual basis. Thus there are \(s_a\in\mathcal O_L\) and
\(\mathcal O_K\)-linear maps \(\varphi_a:\mathcal O_L\to
\mathcal O_K\) such that
\[
                     u=\sum_a s_a\varphi_a(u)
                               \quad(u\in\mathcal O_L).
 \tag{DI.43}
\]
Localization extends each \(\varphi_a\) to \(L\to K\), and
(DI.43) then holds on \(L\).

Take \(z=\sum_j a_jm_j\in W_{\rm good}\), with
\(a_j\in\mathcal O_L\) and \(m_j\in M_K\). Then
\[
 T_{\varphi_a}z=\sum_j\varphi_a(a_j)m_j
                                \in M_K\cap W_{\rm good}.
\]
Membership in \(M_K\) uses its actual \(\mathcal O_K\)-module
law: its defining nonarchimedean norms are at most one, sums satisfy
the ultrametric inequality, and integral scalars have norm at most
one. Equation (DI.43) now writes \(z\) as an
\(\mathcal O_L\)-linear combination of these good integral vectors.
The reverse inclusion in (DI.41) is immediate, proving equality.

For the original \(A=R.FW(i)\), the field
\[
               \mathrm{Holds}.\mathrm{aux}.\mathrm{integral}(i)
\]
gives
\(\mathrm{IsKsForm}\) and \(W_K\)-stability.
\(\mathrm{Pins}.\mathrm{keep}.\mathrm{valKs}\), the residue
character restriction and the group-character equality give
\(\chi(g)\in K_s\): each determinant residue character is the
value of the stored Hecke character on its integral local-unit lift,
and the finite product is the original \(\chi(g)\). Injectivity
of \(\sigma\) identifies the field element before embedding.
\(\mathrm{keep}.\mathrm{levKs}\) gives the values of \(\psi_0\)
and \(e(g)\); inversion preserves the field. Thus every hypothesis
of (DI.41) is verified on the literal data.
% END DI17

% BEGIN DI18
\paragraph{DI18: the original target basis is independent and spans.}
\label{claim:DI18}
For every \(i<m\), the vectors
\(Z.j_W(i)(e_k)\) are complex-linearly independent and span the
whole original \(D.\mathrm{what}_2.M_j(i+1)\).

\emph{Proof.}
Put \(e_i=Y.\mathrm{mj}.e_W(i)\). Its original norm equality is
\(\|e_i x\|^2=n_0\|x\|^2\), with \(n_0>0\); therefore it is
injective. The reading \(Z.\mathrm{jW\_read}\) says exactly
\[
 e_i(Z.j_W(i)y)
   =j_i\left(\sum_k y_k b_{i,k}\right)
   =\mathrm{goodForm}_i(y),
 \tag{DI.44}
\]
where \(b_{i,k}\) is the stored \(L\)-basis of
\(R.WL(F,0,i)\) and \(j_i=(R.FW(i)).j(\mathrm{sg}_0)\).
The base place in DI16 gives \(\mathrm{sg}_0=\sigma\) exactly.

The uniform coefficient-form property
\(\mathrm{Holds}.\mathrm{prim}.\mathrm{unif}\) makes the complex
images of this independent \(L\)-basis independent. By (DI.44)
the original \(Z.j_W(i)e_k\) are independent as well. Their images
span \(R.WS(F,0,i,\mathrm{sg}_0)\): its stored basis is
\(b_C=n_0^{-1/2}j_i(b_L)\), so multiplying by the nonzero
\(\sqrt{n_0}\) changes no span. The original
\(\mathrm{Holds}.\mathrm{aux}.\mathrm{levelW}\) identifies that
space with the full target multiplicity space, and
\(Y.\mathrm{mj}.\mathrm{eW\_range}\) is exactly the range of
\(e_i\). Therefore these images span \(\operatorname{range}e_i\).
Injectivity of \(e_i\) pulls the span equality back to the whole
domain. This includes a zero-dimensional target; the empty basis
then spans its zero space.
% END DI18

% BEGIN DI19
\paragraph{DI19: the exact target integral-lattice tie.}
\label{claim:DI19}
For every \(i<m\),
\[
 \mathrm{OLatTie}\bigl(\sigma,D.\mathrm{what}_2.M_j\mathrm{Lat}(i+1),
                         Z.j_W(i),S_0.M(i)\bigr)
 \tag{DI.45}
\]
holds with the original finite generators and scalar ring.

\emph{Proof.}
Use \(e_i,j_i\) from DI18. The target-lattice reading has the
exact equivalence
\[
 x\in M_j\mathrm{Lat}(i+1)
 \ \Longleftrightarrow\ 
 \exists y\in M_K\cap W_{\rm good},\quad e_i x=j_i y.
 \tag{DI.46}
\]
Every right-hand \(j_i y\) actually lies in the image of \(e_i\).
By DI17, \(y\in M_j\). In its stored basis, write
\(y=\sum_k a_kb_{i,k}\). The definition of the coordinate lattice
\(S_0.M(i)\) says exactly that \(a\) belongs to that lattice.
Equation (DI.44) then supplies the preimage \(Z.j_W(i)a\).
Thus (DI.46) gives an equality of sets of images, not just one
inclusion.

For clarity, here is the generator calculation used to pass to the
literal \(\mathrm{OLatTie}\). If an additive semilinear map \(j\)
satisfies \(j(az)=\sigma(a)j(z)\), the additive subgroup generated
by \(\{\sigma(a)j(z):a\in\mathcal O_L,z\in X\}\) is unchanged
when \(X\) is replaced by its \(\mathcal O_L\)-span. One inclusion
is immediate. For the other, expand a finite linear combination,
distribute \(j\), and multiply the integral coefficients; closure
under addition and negation gives the assertion. Apply this to
(DI.46) and then DI17. The resulting subgroup is exactly
\(e_i(Z.j_W(i)(S_0.M(i)))\).

Let \(\mathcal G\) be the original finite generating set of the
abstract \(M_j\mathrm{Lat}\). Every generator belongs to its
carrier. Conversely every carrier vector is a finite combination
of \(\mathcal G\) over the original coefficient ring \(D.O\).
The map \(D.O\to L\) has integral image by \(D.O_{\rm int}\).
Therefore the subgroup generated over \(\mathcal O_L\) by the
embedded carrier equals the subgroup of all finite
\(\mathcal O_L\)-linear combinations of \(\mathcal G\).
Together with the preceding paragraph, injectivity of \(e_i\)
shows that each generator is in the rational image of \(Z.j_W(i)\),
and that a coordinate vector \(a\) lies in \(S_0.M(i)\) exactly
when \(Z.j_W(i)a\) is such an integral combination of the original
generators. DI18 makes \(Z.j_W(i)\) injective: expand in its
complex-independent coordinate vectors and use injectivity of
\(\sigma\). These are both clauses of (DI.45). No saturation,
topological completion, or replacement lattice is used.
% END DI19

% BEGIN DI20
\paragraph{DI20: the complete first \(\mathrm{DivDef3}\) on the given readers.}
\label{claim:DI20}
The hypotheses of DI1, including the explicit additional
\(\mathrm{RowPairInt}\) supplier (DI.3a), imply
\[
 \begin{split}
 \mathrm{DivDef3}\bigl(&D,A_3,F.P_0,F.\psi_0,F.P_1,F.\psi_1,
        \mathrm{zsOf}(R,F),\mathrm{s0Of}(R,F),\mathrm{s1Of}(R,F),\\
       &X_0,\mathrm{colOf}(R,F,s_{EE}),C.\mathrm{NuC},
                                         \mathrm{idsOf}(c_G,Z,C)\bigr).
 \end{split}
 \tag{DI.47}
\]

\emph{Proof.}
We specify every field. The rank field is \(m=D.R.m\) by its
index type, followed by the two rank equalities in DI16. The group
field is (DI.4), now applied to the two proved cardinality equalities
in DI16. The tuple field uses the original
\(Y.\mathrm{tup0}\), \(Y.\mathrm{muAt}\), DI16's distinguished
place, and injectivity of \(Y.\mathrm{tup}\), which follows by
applying its left inverse \(Y.\mathrm{tupT}\). The equality
\(S_1.\mu=S_0.\mu\) is literal: both definitions use \(R.\mu\).

The sharp and base fields are DI15. The \(\mathrm{mt}\) subrecord
has exactly its two original supplied fields
\(Y.\mathrm{formsL}\) and \(Y.\mathrm{latL}\). The first theta
subrecord \(\mathrm{th}\) has exactly three fields: DI18 gives
\(\mathrm{li}\) and \(\mathrm{span}\), and DI19 gives
\(\mathrm{lat}\). In particular this subrecord has no
\(\mathrm{thTheta}\) field. Finally the column subrecord is exactly
\(C.\mathrm{bij},C.\mathrm{nu}\): the original word bijection and
the full lattice/metric tie under \(C.e_\nu\). Its ambient column
datum is \(\mathrm{colOf}(R,F,s_{EE})\), with that same specified
\(s_{EE}\). The first definition does not inspect a new period
comparison or require a different \(s_{EE}\).

These are all nine fields of \(\mathrm{DivDef3}\). Although the
source's combined theorem also mentions \(Xd\) and \(JM\) for its
second conclusion, the first proof does not use them. The present
proof therefore closes the first-output implication without assuming
\(\mathrm{MjHTie,IdxDefs,BuRational,BtVanish}\), which belong to
the distinct second-output route. It constructs no instance of those
predicates by projecting the first record.
It remains conditional on (DI.3a). In particular it is a proved
sufficient refinement of the original first-output theorem, rather
than a claim that every original premise of that theorem has already
been independently turned into all required analytic inputs.
% END DI20

% BEGIN DI21
\paragraph{DI21: installation on CSR's updated realization and exact remaining scope.}
\label{claim:DI21}
Apply DI20 to the literal
\[
                 R',F',Y',Z',C',s_{EE}
 \tag{DI.48}
\]
of CSR11--CSR13. This supplies the formerly retained first
\(\mathrm{DivDef3}\) input in CSR13, so its actual
\(\mathrm{TwIn}\), \(\mathrm{TwArgs}\), and the subsequent
CSR14--CSR22 finite constructors now require only the explicit
row-pair integrability supplier, together with their other retained
inputs, in place of an unexpanded division-identity premise.

\emph{Proof.}
CSR11 changes only inactive scalar values and proves equality of
both exact row metrics. CSR12 transports \(\mathrm{Pins}\),
\(\mathrm{Holds}\), the same \(\mathrm{Read}\), and the same
\(\mathrm{ReadL}\); its thirteen map identities include all maps
used in DI13, DI15 and DI18--DI19. The same original
\(\mathrm{ReadC}\), chosen before later \(\mathrm{TwArgs}\), is
copied to \(C'\) in CSR13. Its six fields retain its actual
\(w_U,e_\nu,\mathrm{NuC}\), point identity, bijection and metric
tie. Its group tie is the retained original rule (DI.4). Require
(DI.3a) for these same updated data; it is not inferred from their
record equality alone. Every input of DI20 then holds on (DI.48),
with the same geometry, selected columns and period data. No
existential witness in the proof above changes one of those objects.

The result is pointwise and universal in all the remaining indices:
every tuple, both levels, every inner step, every polynomial, every
smooth tensor, every finite level element, and every row coordinate.
It introduces no new asymptotic threshold. Consequently, in CSR22's
original eventual order, retain one common admissible threshold and
one common realization before making the later \(i,u,w\) choices.
The new implication applies at that same realization. In particular
the original \(\mathrm{NEb}\) premise, universal over its already
specified later inputs, is neither specialized to one convenient
input nor reordered behind those choices.

The remaining existence inputs are still substantive. They include
the original compatible \(R\) and readers; geometric realization and
the same columns/periods; the \(\mathrm{Holds}.\mathrm{prim}\)
rational row/lift tests, Hecke and group equations, positive theta
operators and decompositions; the stored full-chain existence in
\(\mathrm{Holds}.\mathrm{aux}.\mathrm{ex}\); its level and integral
data; the original real-group predicates (DI.2); and the exact
row-pair integrability supplier (DI.3a). The current proof does not
derive that supplier from full cusp decay without the missing
global-growth conversion. The new result proves these inputs'
first division compatibility consequence. It does not produce these
inputs, the separate \(\mathrm{ratG/MLat}\) branch, the geometric
carrier/eigensystem/logarithm matching, or the same-realization host
datum. The stronger \(\mathrm{DivRows3}\) remains separate: its
two row-lattice metrics, rational-lift cut, actual \(\mathrm{thTheta}\)
square, index-number identities and column-pairing identity require
their own exact reconstruction and matching.
% END DI21
